\chapter*{Preface}
\addcontentsline{toc}{chapter}{Preface}

\section*{Purpose of This Work}
% Write 1–2 paragraphs that:
% - Explain why you undertook this work (academic, professional, or personal motivation).
% - State the intended contribution: theoretical insight, practical tool, pedagogical resource, etc.
% - Clarify the value for readers (why they should care about this work).

I undertook this work because I was not satisfied with the existing
tools for illustrating scenes with \(3\)D parametric objects in
\TeX{}, such as surfaces. From the perspective of computer graphics
as a whole, the ability to occlusion-order non-intersecting and
non-cyclically overlapping \(0\)--\(2\)-dimensional affine simplices
in \(3\)D---together with the systematic partitioning of simplices to
eliminate intersections and cyclic overlaps---is a foundational
advance. While my own motivation arose from the needs of a
\TeX{} illustrator, I recognized that this framework represents a
paradigm shift in mathematics illustration and, more broadly, a
landmark contribution to computer graphics. My goals were therefore
twofold: to create the kind of illustration tool I had long
envisioned, and to establish priority in implementing this
transformative idea.

I have developed a systematic method for occlusion ordering
\(0\)--\(2\)-dimensional affine simplices in three dimensions.
In addition, I have implemented a framework for eliminating
intersections through partitioning, with the resolution of
cyclic overlaps planned for the next version of the software.
This work matters because it provides the only known approach
for generating scenes with parametric objects that both
intersect coherently and are occlusion-sorted robustly. In
short, this is a rigorous and comprehensive software framework
for handling intersections and occlusion ordering of affine
simplices in three dimensions.

\section*{Context and Background}
% Write 1–2 paragraphs that:
% - Place the work in its institutional or research context (e.g., thesis project, funded research, independent study).
% - If relevant, describe the broader academic or industrial landscape in which the work sits.
% - Briefly acknowledge any special circumstances (e.g., interdisciplinary nature, collaboration across fields).

This book documents the \LaTeX{} package \luatikztdtools{}, which
implements a groundbreaking algorithm for \(3\)D illustrations
composed of \(0\)--\(2\)-dimensional affine simplices. Although the
software is designed for creating illustrations in \LaTeX{}, the
underlying method represents a significant advance for the broader
field of computer graphics.

\section*{Acknowledgments of Guidance and Support}
% Write 1–2 paragraphs that:
% - Express gratitude to supervisors, mentors, or advisors who provided intellectual or technical guidance.
% - Recognize collaborators, peers, or communities that influenced the work.
% - Acknowledge institutional or financial support (scholarships, grants, lab access, software/hardware resources).

I am deeply grateful to the \TeX{} Stack Exchange community, whose
many contributions and generous assistance have been invaluable in
helping me develop my skills as an illustrator. I also acknowledge
the Math Stack Exchange community for several insights that proved
helpful along the way. This work would not have been possible without
the years of learning made available by my peers on \TeX{} Stack
Exchange, nor without the foundational training in programming and
linear algebra that I received from my professors at my local
university.

\section*{Author's Contribution and Perspective}
% Write 1–2 paragraphs that:
% - Clarify your personal role and unique contributions (e.g., system design, proofs, experiments, visualizations).
% - If the work was collaborative, define what you specifically contributed.
% - Optionally share a personal note on your perspective or journey (why this problem matters to you).

I am the principal designer of nearly all of \luatikztdtools{},
with a few specific components assisted by ChatGPT. My two
primary original contributions are the occlusion-ordering
algorithm and the simplex-partitioning algorithm. The portion
produced with ChatGPT's assistance concerns the recursive
traversal of segments that applies my algorithms. Apart from
this limited assistance, the work was carried out independently,
with valuable input and feedback from the \TeX{} Stack Exchange
community.

\section*{Structure of the Document}
% Write 1 short paragraph that:
% - Give a roadmap of the thesis/book: briefly describe what each chapter covers.
% - Keep it concise (one sentence per chapter).
% - Help orient the reader before they dive into the main content.

This book is organized into five chapters. Chapter~1 presents the
problem statement, defining and framing the central challenge.
Chapter~2 reviews the relevant literature. Chapter~3 describes the
methodology, including the algorithms developed. Chapter~4 provides
validation of the approach, and Chapter~5 presents the results and
analysis.